\documentclass[11pt,oneside]{book}

\usepackage[a4paper,margin=2.7cm,headheight=28pt]{geometry}
\usepackage{amsmath,amssymb}
\usepackage{array,longtable}
\usepackage{graphicx}
\usepackage{tikz}
\usepackage{pgfplots}
\usepackage{enumitem}
\usepackage{fancyhdr}
\usepackage[hidelinks]{hyperref}
\usepackage{xcolor}
\usepackage{indentfirst}
\usepackage{siunitx} % For consistent units formatting
\usepackage{enumitem}

\pgfplotsset{compat=1.18}
\linespread{1.15}
\setlength{\parindent}{2em}
\setlength{\parskip}{0.3em}
\setlength{\headsep}{18pt}
\setlength{\skip\footins}{10pt}

\hypersetup{
  pdftitle={The Principle of Relativity},
  pdfauthor={***}
}

\pagestyle{fancy}
\fancyhf{}
\fancyhead[L]{\parbox[b]{0.85\textwidth}{\raggedright\nouppercase{\leftmark}}}
\fancyhead[R]{\thepage}
\renewcommand{\headrulewidth}{0.4pt}
\renewcommand{\footrulewidth}{0pt}

% Keep section numbers local to each paper instead of chapter.section.
\renewcommand{\thesection}{\arabic{section}}
\renewcommand{\thesubsection}{\thesection.\arabic{subsection}}
\renewcommand{\theHsection}{\thechapter.\arabic{section}}
\renewcommand{\theHsubsection}{\thechapter.\arabic{section}.\arabic{subsection}}
\renewcommand{\theequation}{\arabic{equation}}

\newtheorem{theorem}{Theorem}[chapter]
\newtheorem{proposition}[theorem]{Proposition}
\newtheorem{lemma}[theorem]{Lemma}
\newtheorem{definition}[theorem]{Definition}
\newtheorem{remark}[theorem]{Remark}
\newenvironment{proof}{\par\noindent\textbf{Proof.} }{\hfill$\square$\par}

\newcommand{\papertitle}[1]{%
  {\centering\Huge\bfseries #1\par}
}

\newcommand{\papersubtitle}[1]{%
  \vspace{0.6em}
  {\centering\large\itshape #1\par}
}

\newcommand{\paperauthor}[1]{%
  \vspace{1.2em}
  {\centering\large #1\par}
}

% \newcommand{\paperinfo}[2]{%
%    \vspace{0.8em}
%   {\centering\small
%   \begin{tabular}{c}
%   \textbf{Brief Introduction}\quad #1\\
%   \textbf{Publish Time}\quad #2
%   \end{tabular}\par}%
% }

\newcommand{\paperheader}[5]{%
    \refstepcounter{chapter}
    \chapter*{}
    \thispagestyle{plain}
    \papertitle{#1}
    \paperauthor{#2}
    \papersubtitle{#3}
    %\paperinfo{#4}{#5}
    \vspace{1.3em}
    \noindent\rule{\textwidth}{0.6pt}
    \vspace{1.2em}
    \addcontentsline{toc}{chapter}{#1}
    \markboth{#1}{#1}
}

\newcommand{\collectiontitle}{\Huge{THE\\PRINCIPLE\\OF\\RELATIVITY}}
\newcommand{\collectionsubtitle}{\huge{A Collection of Original Papers on the Special and General Theories of Relativity}}
\newcommand{\collectioneditor}{}
\newcommand{\collectiondate}{2026 B.C.}

\begin{document}
\color{black}

\frontmatter
\begin{titlepage}
    \centering
    \vspace*{3.5cm}
    {\fontsize{24}{30}\selectfont\bfseries\collectiontitle\par}
    \vspace{1.2cm}
    {\Large\collectionsubtitle\par}
    \vspace{2.2cm}
    {\large\collectioneditor\par}
    \vfill
    {\large\collectiondate\par}
    \vspace*{1.5cm}
\end{titlepage}

\newgeometry{left=4.5cm,right=4.5cm,top=3cm,bottom=3cm,headheight=28pt}
\chapter*{Preface}
\addcontentsline{toc}{chapter}{Preface}
\Large
This collection of papers has been reproduced from the book \textit{The Principle of Relativity} (ISBN: 978-0-486-60081-9), originally published by Dover Publications in 1952, translated by W. Perrett and G. B. Jeffery, notes by A. Sommerfeld. What the editor has done is just to make minor typographical adjustments to the format and layout and add some comments.

The papers herein are arranged in chronological order, and the original text has been preserved as faithfully as possible. The purpose of this volume is to provide a comprehensive and accessible introduction to the theory of relativity, allowing readers to engage directly with the foundational works of the great physicists who shaped this profound branch of physics.

In addition, we have provided references to the original sources of these papers. We strongly encourage interested readers to consult these primary documents to cultivate a deeper and more authentic understanding of the principles of relativity.
\par\normalsize

\restoregeometry


\tableofcontents

\mainmatter
\paperheader
  {Michelson's Interference Experiment}
  {H. A. Lorentz}
  {Translated from "Versuch einer Theorie der elektrischen und optischen Erscheinungen in bewegten Körpern," Leidem, 1895, \S\S\,89-92.}


\section*{1}
As Maxwell first remarked and as follows from a very simple calculation, the time required by a ray of light to travel frome a point A to a point B and back to A must vary when the two points together undergo a displacement without carrying the ether with them. The differnce is, certainly, a magnitude of second order; but it is sufficiently great to be detected by a sensitive interference methiod.

The experiment was carried out by Michelson in 1881.\footnote{Michelson, American Journal of Science, 22, 1881, p. 120.} His apparatus, a kind of interferometer, had two hrizontal arms, P and Q, of equal length and at right angles one to the other. Of the two mutually interfering rays of light the one passed along the arm P and back, the other along the arm Q and back. The whole instrument, including the source of light and the arrangement for taking observations, could be recolved about a vertical axis; and those two positions come especially under consideration in which the arm P or the arm Q lay as nearly as possible in the direction of the Earth's motion. On the basis of Fresnel's theory it was anticipated that when the apparatus was revolved from one of these \textit{principal positions} into the other there would be a displacement of the interference fringes.

But of such a displacement - for the sake if brevity we will call it the Maxwell displacement - conditioned by the change in the times of propagation, no trace was discoverd, and accordingly Michelson thought himself justified in concluding that while the Earth is moving, the ether does not remain at rest. The correctness of this inference was soon brought into question, for by an oversight Michelson had taken the change in the phase differnce, which was to be expected in accordance with the theory, at twice its proper value. If we make the necessary correction, we arrive at displacements no greater than might be masked by errors of observation.

Subsequently Michelson \footnote{Michelson and Morley, American Journal of Science, 34, 1887, p. 333; Phil. Mag., 24, 1887, p. 449.} took up the investigation anew in collaboration with Morley, enhancing the delicacy of the experiment by causing each pencil to be reflected to and from between a number of mirrors, thereby obtaining the same advantage as if the arms of the earlier apparatus had been considerably lengthened. The mirrors wre mounted on a massive stone disc, floating on mercury, and therefore easily revolved. Each pencil now had to travel a total distance of 22 meters, and on Fresnel'stheory the displacement to be expected in passing from the one principal position to the other would be 0.4 of the distance between the interference fringes. Nevertheless the rotation produced displacements not exceeding 0.02 of this distance, and these might well be ascribed to errors of observation.

Now, does this result entitle us to assume that the ether takes part in the motion of the Earth, and therefore that the theory of aberration given by Stokes is incorrect one? The difficulties which this theory encounters in explaining aberration seem too great for me to share this opinion, and I would rather try to remove the contradiction between Frensel's theory and Michelson's result. An hypothesis which I brought forward some time ago, \footnote{Lorentz, Zittingsverslagen der Akad. v. Wet. te Amsterdam, 1892-93, p.74.} and which, as I subsequently learned, has also occurred to Fitzgerald, \footnote{As Fitzgerald kindly tells me, he has for a long time dealt with his hypothesis in his lectures. The only published reference which I can find to the hypothesis is by Lodge, "Aberration Problems," Phil. Trans. R.S., 184A, 1893.} enables us to do this. The next paragraph will set out this hypothesis.


\section*{2}
To simplify matters we will assume that we are working with apparatus as employed in the first experiment, and that in the one principal position the arm P lies exactly in the direction of the Earth's motion. Let $v$ be the velocity of this motion, $L$ the length of either arm, and hence $2L$ the path traversed by the rays of light. Accordin to the theory, \footnote{Cf. Lorentz. Arch. Néerl., 2, 1887, pp.168-176.}

\paperheader
    {Electromagnetic Phenomena in a System Moving with Any Velocity Less Than That of Light}
    {H. A. Lorentz}
    {...}


\section{Introduction}
hello world
\paperheader
	{On the Electrodynamics of Moving Bodies}
	{A. Einstein}
	{Translated from "Zur Elektrodynamik bewegter Körper," Annalen der Physik, 17, 1905.}


\section*{}
It is known that Maxwell's electrodynamics - as usually understood at the present time - when applied to moving bodies, leads to asymmetries which do not appear to be inherent in the phenomena. Take, for example, the reciprocal electrodynamic action of a magnet and a conductor. The observable phenomenon here depends only on the relative motion of the conductor and the magnet, whereas the customary view draws a sharp distinction between the two cases in which either the one or the other of these bodies is in motion. For if the magnet is in motion and the conductor at rest, there arises in the neighbourhood of the magnet an electric field with a certain definite energy, producing a current at the places where parts of the conductor are situated. But if the magnet is stationary and the conductor in motion, no electric field arises in the neighbourhood of the magnet. In the conductor, however, we find an electromotive force, to which in itself there is no corresponding energy, but which gives rise - assuming the equality of the relative motion in  the two cases discussed - to electric currents of the same path and intensity as those produced by the electric forces in the former case.

Examples of this sort, together with the unsuccessful attempts to discover any motion of the earth relative to the "light medium," suggest that the phenomena of electrodynamics as well as of mechanics possess no properties corresponding to the idea of absolute rest. They suggest rather that, as has already been shown to the first order of small quantities, the same laws of electrodynamics and optics will be valid for all frames of reference for which the equations of mechanics hold good.\footnote{The preceding memoir by Lorentz was not at this time known to the author.} We will raise this conjecture (the purport of which will hereafter be called the "Principle of Relativity") to the status of a postulate, and also introduce another postulate, which is only apparently irreconcilable with the former, namely, that light is always propagated in empty space with a definite velocity $c$ which is independent of the state of motion of the emitting body. These two postulates suffice for the attainment of a simple and consistent theory of the electrodynamics of moving bodies based on Maxwell's theory for stationary bodies. The introduction of a "luminiferous ether" will prove to be superfluous inasmuch as the view here to be developed will not require an "absolutely stationary space" provided with special properties, nor assign a velocity-vector to a point of the empty space in which electromagnetic processes take place.

The theory to be developed is based - like all electrodynamics - on the kinematics of the rigid body, since the assertions of any such theory have to do with the relationships between rigid bodies (systems of co-ordinates), clocks, and electromagnetic processes. Insufficient consideration of this circumstance lies at the root of the difficulties which the electrodynamics of moving bodies at present encounters.


\vspace{2em}
\begin{center}
    \Large\textbf{I. KINEMATICAL PART}
\end{center}

\section{Definition of Simultaneity}
Let us take a system of co-ordinates in which the equations of Newtonian mechanics hold good.\footnote{i.e. to the first approximation.} In order to render our presentation more precise and to distinguish this system of co-ordinates verbally from others which will be introduced hereafter, we call it the "stationary system."

If a material point is at rest relatively to this system of co-ordinates, its position can be defined relatively thereto by the employment of rigid standards of measurement and the methods of Euclidean geometry, and can be expressed in Cartesian co-ordinates.

If we wish to describe the \textit{motion} of a material point, we give the values of its co-ordinates as functions of the time. Now we must bearcarefully in mind that a mathematical description of this kind has no physical meaning unless we are quite clear as to what we understand by "time." We have to take into account that all our judgments in which time plays a part are always judgments of \textit{simultaneous events}. If, for instance, I say, "That train arrives here at 7 o'clock," I mean something like this: "The pointing of the small hand of my watch to 7 and the arrival of the train are simultaneous events."\footnote{We shall not here discuss the inexactitude which lurks in the concept of simultaneity of events at approximately the same place, which can only be removed by an abstraction.}

It might appear possible to overcome all the difficluties attending the definition of "time" by substituting "the position of the small hand of my watch" for "time". And in fact such a definition is satisfactory when we are concerned with defining a time exclusively for the place where the watch is located; but it is no longer satisfactory when we have to connect in time series of events occurring at different places, or - what comes to the same thing - to evaluate the times of events occurring at places remote from the watch.

We might, of course, content ourselves with time values determined by an observer stationed together with the watch at the origin of the co-ordinates, and co-ordinatings the corresponding positions of the hands with light signals, given out by every event to be timed, and reaching him through empty space. But this co-ordination has the disadvantage that it is not independent of the standpoint of the observer with the watch or clock, as we know from experience. We arrive at a much more practical determination along the following line of thought.

If at the point $A$ of space there is a clock, an observer at $A$ can determine the time values of events in the immediate proximity of $A$ by finding the positions of the hands which are simultaneous with these events. If there is at the point $B$ of space another clock in all respects resembling the one at $A$, it is possible for an observer at $B$ to determine the time values of events in the immediate neibourhood of $B$. But it is not possible without further assumption to compare, in respect of time, an event at $A$ with an event at $B$. We have so far defined only an "A time" and a "B time." We have not defined a common "time" for $A$ and $B$, for the latter cannot be defined at all unless we establish \textit{by definition} that the "time" it requires to travel from $A$ to $B$ equals the "time" it requires to travel from $B$ to $A$. Let a ray of light start at the "A time" $t_A$ from $A$ towards $B$, let it at the "B time" $t_B$ be reflected at $B$ in the direction of $A$, and arrive again at $A$ at the "A time" $t'_A$.  

In accordance with the definition the two clocks synchronize if
$$ t_B - t_A = t'_A - t_B. $$

We assume that this definition of synchronism is free from contradictions, and possible for any number of points; and that the following relations are universally valid: -
\begin{enumerate}
	\item If the clock at $B$ synchronizes with the clock at $A$, then the clock at $A$ synchronizes with the clock at $B$.
	\item If the clock at $A$ synchronizes with the clock at $B$ and also with the clock at $C$, the clocks at $B$ and $C$ also synchronize with each other.
\end{enumerate}

Thus with the help of certain imaginary physical experiments we have settled what is to be understood by synchronous stationary clocks located at different places, and have evidently obtained a definition of "simultaneous", or "synchronous", and of "time." The "time" of an event is that which is given simultaneously with the event by a stationary clock located at the place of the event, this clock being synchronous, and indeed synchronous for all time determinations, with a specified stationary clock.

In agreement with experience we further assume that the quantity
$$ \frac{2AB}{t'_A - t_A} = c ,$$
to be a universal constant - the velocity of light in empty space.

It is essential to have time defined by means of stationary clocks in the stationary system, and the time now defined being appropriate to the stationary system we call it "the time of the stationary system."


\section{On the Relativity of Lengths and Times}
The following reflexions are based on the principle of relativity and on the principle of the constancy of the velocity of light. These two principles we define as follows: -
\begin{enumerate}%[leftmargin=2.2em]
	\item The laws by which the states of physical systems undergo change are not affected, whether these changes of state be referred to the one or the other of two systems of co-ordinates in uniform translatory motion.
	\item Any ray of light moves in the "stationary" system of co-ordinates with the determined velocity $c$, whether the ray be emitted by a stationary or by a moving body. Hence $$\text{velocity} = \frac{\text{light path}}{\text{time interval}}$$ where time interval is to be taken in the sense of the definition in \S 1.
\end{enumerate}

Let there be given a stationary rigid rod; and let its length be $l$ as measured by a measuring-rod which is also stationary. We now imagine the axis of the rod lying along the the axis of $x$ in the direction of increasing $x$ is then imparted to the rod. We now inquire as to the length of the moving rod, and imagine its length to be ascertained by the following two operations: -
\begin{enumerate}[label=(\alph*).]
	\item The observer moves together with the given measuring-rod and the rod to be measured, and measures the rod directly by superposing the measuring-rod, in just the same way as if all three were at rest. \label{item:measurement1}
	\item By means of stationary clocks set up in the stationary system and synchronizing in accordance with \S 1, the observer ascertains at what points of the stationary system the two ends of the rod to be measured are located at a definite time. The distance between these two points, measured by the measuring-rod already employed, which in this case is at rest, is also a length which may be designated "the length of the rod." \label{item:measurement2}
\end{enumerate}

In accordance with the principle of relativity the length to be discovered by the operation \ref{item:measurement1} - we will call it "the length of the rod in the moving system" - must be equal the length $l$ of the rod in the stationary rod.

The length to be discovered by the operation \ref{item:measurement2} we will call "the length of the (moving) rod in the stationary system." This we shall determine on the basis of our two principles, and we shall find that it differs from $l$.

Current kinematics tacitly assumes that the lengths determined by those two operations are precisely equal, or in other words, that a moving rigid body at the epoch $t$ may in geometrical respects be perfectly represented by \textit{the same} body at \textit{rest} in a definite position.

We imagine further that at the two ends $A$ and $B$ of the rod, clocks are placed which synchronize with the clocks of the stationary system, that is to say that their indications correspond at any instant to the "time of the stationary system" at the places where they happen to be. These clocks are therefore "synchronous in the stationary system."

We imagine further that with each clock there is a moving observer, and that these observers apply to both clocks the criterion established in \S 1 for the synchronization of two clocks. Let a ray of light depart from $A$ at the time\footnote{"Time" here denotes "time of the stationary system" and also "position of hands of the moving clock situated at the place under discussion."} $t_A$, let it be reflected at $B$ at the time $t_B$, and reach $A$ again at the time $t'_A$. Taking into consideration the principle of constancy of the velocity of light we find out that $$ t_B - t_A = \frac{r_{AB}}{c-v} \quad\text{and}\quad t'_A - t_B = \frac{r_{AB}}{c+v} $$ where $r_{AB}$ donates the length of the moving rod - measured in the stationary system. Observers moving with the moving rod would thus find that the two clocks are not synchronous, while observers in the stationary system would declare the clocks to be synchronous.

So we see that we cannot attach any \textit{absolute} signification to the concept of simultaneity, but that two events which, viewed from a system of co-ordinates, are simultaneous, can no longer be looked upon as simultaneous events when envisaged from a system which is in motion relatively to that system.


\section{Theory of the Transformation of Co-ordinates and Times from a Stationary System to another System in Uniform Motion of Translation Relative to the Former}
Let us in "stationary" space take two systems of co-ordinates, i.e. two systems, each of three rigid material lines, perpendicular to one another, and issuing from a point. Let the axes of $X$ of the two systems coincide, and their axes of $Y$ and $Z$ respectively be parallel. Let each system be provided with a rigid measuring-rod, and likewise all the clocks of the two systems, be in all respects alike.

Now to the origin of one of the two systems ($k$) let a constant velocity $v$ be imparted in the direction of the increasing $x$ of the other stationary system ($K$), and let this velocity be communicated to the axes of the co-odinates, the relevant measuring-rod, and the clocks. To any time of the stationary system $K$ there then will correspond a definite position of the axes of the moving system, and from reasons of symmetry we are entitled to assume that the motion of $k$ may be such that the axes of the moving system are at the time $t$ (this "$t$" always denotes a time of the stationary system) parallel to the axes of the stationary system.

We now imagine space to be measured from the stationary system $K$ by means of the stationary measuring-rod, and also from the moving system $k$ by means of the measuring-rod moving with it; and that we thus obtain the co-ordinates $x$, $y$, $z$, and $\xi$, $\eta$, $\zeta$ respectively. Further, let the time $t$ of the stationary system be determined for all points thereof at which there are clocks by means of light signals in the manner indicated in \S 1; similarly let the time $\tau$ of the moving system be determined for all points of the moving system at which there are clocks at rest relatively to that system by applying the method, given in \S 1, of light signals between the points at which the latter clocks are located.

To any system of values $x$, $y$, $z$, $t$, which completely defines the place and time of an event in the stationary system, there belongs a system of values $\xi$, $\eta$, $\zeta$, $\tau$, determining that event relatively to the system $k$, and our task is now to find the system of equations connecting these quantities.

In the first place it is clear that the equations must be \textit{linear} on account of the properties of homogeneity which we attribute to space and time.

If we place $x' = x - vt$, it is clear that a point at rest in the system $k$ must have a system of values $x'$, $y$, $z$, independent of time. We first define $\tau$ as a function of $x'$, $y$, $z$, and $t$. To do this we have to express in equations that $\tau$ is nothing else than the summary of the data of clocks at rest in system $k$, which have been synchronized according to the rule given in \S 1.

From the origin of the system $k$ let a ray be emitted at the time $\tau_0$ along the $X$-axis to $x'$, and at time $\tau_1$ be reflected thence to the origin of the co-odinates, arriving there at the time $\tau_2$; we then must have $\frac{1}{2}(\tau_0 + \tau_2) = \tau_1$, or, by inserting the arguments of the functions $\tau$ and applying the principle of the constancy of the velocity of light in the stationary system : - $$ \frac{1}{2}\left[ \tau(0,0,0,t) + \tau\left(0,0,0,t+\frac{x'}{c-v}+\frac{x'}{c+v}\right) \right] = \tau\left( x',0,0, t+\frac{x'}{c-v} \right) .$$ Hence, if $x'$ be chosen infinitesimally small, $$ \frac{1}{2}\left( \frac{1}{c-v} + \frac{1}{c+v} \right) \frac{\partial\tau}{\partial t} = \frac{\partial \tau}{\partial x'} + \frac{1}{c-v}\frac{\partial\tau}{\partial t} ,$$ or $$ \frac{\partial\tau}{\partial x'} = \frac{v}{c^2 - v^2} \frac{\partial\tau}{\partial t} .$$

It is to be noted that instead of the origin of the co-odinates we might have chosen any other point for the point of origin of the ray, and the equation just obtained is therefore valid for all values of $x'$, $y$, $z$.

An analogous consideration - applied to the axes of $Y$ and $Z$ - it being borne in mind that light is always propagated along these axes, when viewed from the stationary system, with the velocity $\sqrt{c^2 - v^2}$, gives us $$ \frac{\partial\tau}{\partial y} = 0,\quad \frac{\partial\tau}{\partial z} = 0 .$$ Since $\tau$ is a \textit{linear} function, it follows from these equations that $$ \tau = a\left( t - \frac{v}{c^2 - v^2} x' \right) $$ where $a$ is a function $\phi(v)$ at present unkown, and where for brevity it is assumed that at the origin of $k$, $\tau = 0$, when $t = 0$.

With the help of this result we easily determine the quantities $\xi$, $\eta$, $\zeta$ by expressing in equations that light (as required by the principle of the constancy of the velocity of light, in combination with the principle of relativity) is also propagated with velocity $c$ when measured in the moving system. For a ray of light emitted at the time $\tau = 0$ in the direction of the increasing $\xi$ $$ \xi = c\tau \quad\text{or}\quad \xi = ac\left( t - \frac{v}{c^2 - v^2} x' \right) .$$ But the ray moves relatively to the initial point of $k$, when measured in the stationary system, with the velocity $c-v$, so that $$ \frac{x'}{c-v} = t .$$ If we insert this value of $t$ in the equation for $\xi$, we obtain $$ \xi = a\frac{c^2}{c^2 - v^2}x' .$$ In an analogous manner we find, by considering rays moving along the two other axes, that $$ \eta = a\tau = ac\left( t - \frac{v}{c^2 - v^2} x' \right) $$ when $$ \frac{y}{\sqrt{c^2 - v^2}} = t ,\quad x' = 0 .$$ Thus $$ \eta = a\frac{c}{\sqrt{c^2 - v^2}}y \quad\text{and}\quad \zeta = a\frac{c}{\sqrt{c^2 - v^2}}z .$$ 

Substituting for $x'$ its value, we obtain
\begin{align*}
	\tau &= \varphi(v)\beta(t - vx/c^2), \\
	\xi &= \varphi(v)\beta(x - vt), \\
	\eta &= \varphi(v)y, \\
	\zeta &= \varphi(v)z,
\end{align*}
where $$\beta = \frac{1}{\sqrt{1 - v^2/c^2}} ,$$ and $\varphi$ is an as yet unknown function of $v$. If no assumption whatever be made as to the initial position of the moving system and as to the zero point of $\tau$, an additive constant is to be placed on the right side of each of these equations.

We now have to prove that any ray of light, measured in the moving system, is propagated with the velocity $c$, if, as we have assumed, this is the case in the stationary system; for we have not as yet furnished the proof that the principle of the constancy of the velocity of light is compatible with the principle of relativity.

At the time $t = \tau = 0$, when the origin of the co-ordinates is common to the two systems, let a spherical wave be emitted therefrom, and be propagated with the velocity $c$ in system $K$. If $(x, y, z)$ be a point just attained by this wave, then $$ x^2 + y^2 + z^2 = c^2 t^2 .$$

Transforming this equation with the aid of our equations of transformation we obtain after a simple calculation $$ \xi^2 + \eta^2 + \zeta^2 = c^2 \tau^2 .$$

The wave under consideration is therefore no less a spherical wave with velocity of propagation $c$ when viewed in the moving system. This shows that our two fundamental principles are compatible.\footnote{The equations of the Lorentz transformation may be more simply deduced directly from the condition that in virtue of those equations the relation $x^2 + y^2 + z^2 = c^2 t^2$ shall have as its consequence the second relation $\xi^2 + \eta^2 + \zeta^2 = c^2 \tau^2$.}

In the equations of transformation which have been developed there enters an unknown function $\varphi$ of $v$, which we will now determine.

For this purpose we introduce a third system of co-ordinates $K'$, which relatively to the system $k$ is in a state of parallel translatory motion parallel to the axis of $X$, such that the origin of co-ordinates of system $k$ moves with velocity $-v$ on the axis of $X$. At the time $t = 0$ let all three origins coincide, and when $t=x=y=z=0$ let the time $t'$ of the system $K'$ be zero. We call the co-ordinates, measured in the system $K'$, $x'$, $y'$, $z'$, and by twofold application of our equations of transformation we obtain
\begin{alignat*}{2}
	t' &= \varphi(-v)\beta(-v)(\tau + v\xi/c^2) &&= \varphi(v)\varphi(-v)t, \\
	x' &= \varphi(-v)\beta(-v)(\xi+v\tau) &&= \varphi(v)\varphi(-v)x, \\
	y' &= \varphi(-v)\eta &&= \varphi(v)\varphi(-v)y, \\
	z' &= \varphi(-v)\zeta &&= \varphi(v)\varphi(-v)z.
\end{alignat*}

Since the relations between $x'$, $y'$, $z'$ and $x$, $y$, $z$ do not contain the time $t$, the systems $K$ and $K'$ are at rest with respect to one another, and it is clear that the transformation from $K$ to $K'$ must be the identical transformation. Thus $$ \varphi(v)\varphi(-v) = 1 .$$ We now inquire into the signification of $\varphi(v)$. We give our attention to that part of the axis of $Y$ of system $k$ which lies between $\xi = 0$, $\eta = 0$, $\zeta = 0$ and $\xi = 0$, $\eta = l$, $\zeta = 0$. This part of the axis of $Y$ is a rod moving perpendicularly to its axis with velocity $v$ relatively to system $K$. Its ends possess in $K$ the co-ordinates $$ x_1 = vt,\quad y_1 = \frac{l}{\varphi(v)},\quad z_1 = 0 $$ and $$ x_2 = vt,\quad y_2 = 0,\quad z_2 = 0 .$$ The length of the rod measured in $K$ is therefore $ l/\varphi(v) $; and this gives us the meaning of the function $\varphi$. From reasons of symmetry it is now evident that the length of a given rod moving perpendicularly to its axis, measured in the stationary system, must depend only on the velocity and not on the direction and the sense of the motion. The length of the moving rod measured in the stationary system does not change, therefore, if $v$ and $-v$ are interchanged. Hence follows that $l/\varphi(v) = l/\varphi(-v)$, or $$\varphi(v) = \varphi(-v) .$$ It follows from this relation and the one previously found that $\varphi(v) = 1$, so that the transformation equations which have been found become
\begin{align*}
	\tau &= \beta(t - vx/c^2), \\
	\xi &= \beta(x - vt), \\
	\eta &= y, \\
	\zeta &= z.
\end{align*}
where $$\beta = 1/\sqrt{1-v^2/c^2}.$$


\section{Physical Meaning of the Equations Obtained in Respect to Moving Rigid Bodies and Moving Clocks}




\section{The Composition of Velocities}




\vspace{2em}
\begin{center}
    \Large\textbf{II. ELECTRODYNAMICAL PART}
\end{center}

\section{Transformation of the Maxwell-Hertz Equations for Empty Space. On the Nature of the Electromotive Forces Occurring in a Magnetic Field During Motion}




\section{Theory of Doppler's Principle and of Aberration}




\section{Transformation of the Energy of Light Rays. Theory of the Pressure of Radiation Exerted on Perfect Reflectors}




\section{Transformation of the Maxwell-Hertz Equations when Convection-Currents are Taken into Account}




\section{Dynamics of the Slowly Accelerated Electron}

\paperheader
	{Does the Inertia of a Body Depend Upon Its Energy Content?}
	{}
	{}


\section{Introduction}
hello world
\paperheader
	{Space and Time}
	{}
	{}

\section{Introduction}
hello world
\paperheader
	{On the Influence of Gravitation on the Propagation of Light}
	{}
	{}

\section{Introduction}
hello world
\paperheader
	{The Foundation of the General Theory of Relativity}
	{}
	{}


\section{Introduction}
hello world
\paperheader
	{Hamiton's Principle and the General Theory of Relativity}
	{}
	{}


\section{Introduction}
hello world
\paperheader
	{Cosmological Considerations on the General Theory of Relativity}
	{}
	{}


\section{Introduction}
hello world
\paperheader
	{Do Gravitational Fields Play an Essential Part in the Structure of the Elementary Particles of Matter?}
	{}
	{}

\section{Introduction}
hello world
\paperheader
	{Glavitation and Electricity}
	{}
	{}


\section{Introduction}
hello world

\backmatter


% \chapter*{Typesetting Notes}
% \addcontentsline{toc}{chapter}{Typesetting Notes}
% To extend this collection, use the following pattern:
% \begin{enumerate}[leftmargin=2.2em]
%   \item Call \verb|\paperheader{Title}{Author}{Subtitle}{Brief introduction}{Publish time}| to create the opening block of each paper.
%   \item Organize the paper body with \verb|\section| and \verb|\subsection|.
%   \item Typeset mathematics with environments such as \verb|equation| and \verb|align|, and keep tables visually restrained and consistent.
%   \item Produce simple diagrams directly with \verb|TikZ| or \verb|PGFPlots| when external image files are unnecessary.
% \end{enumerate}

\end{document}
